% Chunk: Chapter 1 backmatter---Bibliography and References
% Source: Weinberg QFT Vol. I, physical PDF pp. 62--71 (printed pp. 39--48)
% Figures: none.
% Uncertainties: none.

\chapterbackmatter{Bibliography}

\begin{itemize}
\renewcommand{\labelitemi}{$\square$}
\item S. Aramaki, `Development of the Renormalization Theory in Quantum
Electrodynamics,' \emph{Historia Scientiarum} \textbf{36}, 97 (1989); \emph{ibid.}
\textbf{37}, 91 (1989). [Section 1.3.]

\item R. T. Beyer, ed., \emph{Foundations of Nuclear Physics} (Dover
Publications, Inc., New York, 1949). [Section 1.2.]

\item L. Brown, `Yukawa's Prediction of the Meson,' \emph{Centauros}
\textbf{25}, 71 (1981). [Section 1.2.]

\item L. M. Brown and L. Hoddeson, eds., \emph{The Birth of Particle Physics}
(Cambridge University Press, Cambridge, 1983). [Sections 1.1, 1.2, 1.3.]

\item T. Y. Cao and S. S. Schweber, `The Conceptual Foundations and the
Philosophical Aspects of Renormalization Theory,' \emph{Synthese} \textbf{97},
33 (1993). [Section 1.3.]

\item P. A. M. Dirac, \emph{The Development of Quantum Theory} (Gordon and
Breach Science Publishers, New York, 1971). [Section 1.1.]

\item E. Fermi, `Quantum Theory of Radiation,' \emph{Rev. Mod. Phys.}
\textbf{4}, 87 (1932). [Sections 1.2 and 1.3.]

\item G. Gamow, \emph{Thirty Years that Shook Physics} (Doubleday and Co.,
Garden City, New York, 1966). [Section 1.1.]

\item M. Jammer, \emph{The Conceptual Development of Quantum Mechanics}
(McGraw-Hill Book Co., New York, 1966). [Section 1.1.]

\item J. Mehra, `The Golden Age of Theoretical Physics: P.A.M. Dirac's
Scientific Work from 1924 to 1933,' in \emph{Aspects of Quantum Theory}, ed.
by A. Salam and E. P. Wigner, (Cambridge University Press, Cambridge, 1972).
[Section 1.1.]

\item A. I. Miller, \emph{Early Quantum Electrodynamics---A Source Book}
(Cambridge University Press, Cambridge, UK, 1994). [Sections 1.1, 1.2, 1.3.]

\item A. Pais, \emph{Inward Bound} (Clarendon Press, Oxford, 1986).
[Sections 1.1, 1.2, 1.3.]

\item S. S. Schweber, `Feynman and the Visualization of Space-Time
Processes,' \emph{Rev. Mod. Phys.} \textbf{58}, 449 (1986). [Section 1.3.]

\item S. S. Schweber, `Some Chapters for a History of Quantum Field Theory:
1938--1952,' in \emph{Relativity, Groups, and Topology II}, ed. by B. S.
DeWitt and R. Stora (North-Holland, Amsterdam, 1984). [Sections 1.1, 1.2,
1.3.]

\item S. S. Schweber, `A Short History of Shelter Island I,' in \emph{Shelter
Island II}, ed. by R. Jackiw, S. Weinberg, and E. Witten (MIT Press,
Cambridge, MA, 1985). [Section 1.3.]

\item S. S. Schweber, \emph{QED and the Men Who Made It: Dyson, Feynman,
Schwinger, and Tomonaga} (Princeton University Press, Princeton, 1994).
[Section 1.1, 1.2, 1.3.]

\item J. Schwinger, ed., \emph{Selected Papers in Quantum Electrodynamics}
(Dover Publications Inc., New York, 1958). [Sections 1.2 and 1.3.]

\item S.-I. Tomonaga, in \emph{The Physicist's Conception of Nature}
(Reidel, Dordrecht, 1973). [Sections 1.2 and 1.3.]

\item S. Weinberg, `The Search for Unity: Notes for a History of Quantum
Field Theory,' \emph{Daedalus}, Fall 1977. [Sections 1.1, 1.2, 1.3.]

\item V. F. Weisskopf, `Growing Up with Field Theory: The Development of
Quantum Electrodynamics in Half a Century,' 1979 Bernard Gregory Lecture at
CERN, published in L. Brown and L. Hoddeson, \emph{op. cit.}. [Sections 1.1,
1.2, 1.3.]

\item G. Wentzel, `Quantum Theory of Fields (Until 1947)' in \emph{Theoretical
Physics in the Twentieth Century}, ed. by M. Fierz and V. F. Weisskopf
(Interscience Publishers Inc., New York, 1960). [Sections 1.2 and 1.3.]

\item E. Whittaker, \emph{A History of the Theories of Aether and Electricity}
(Humanities Press, New York, 1973). [Section 1.1.]
\end{itemize}

\chapterbackmatter{References}

\begin{enumerate}
\item\label{ref:1.1} L. de Broglie, \emph{Comptes Rendus} \textbf{177}, 507, 548, 630
(1923); \emph{Nature} \textbf{112}, 540 (1923); Th\`ese de doctorat (Masson et
Cie, Paris, 1924); \emph{Annales de Physique} \textbf{3}, 22 (1925)
[reprinted in English in \emph{Wave Mechanics}, ed. by G. Ludwig, (Pergamon
Press, New York, 1968)]; \emph{Phil. Mag.} \textbf{47}, 446 (1924).

\item\label{ref:1.2} W. Elsasser, \emph{Naturwiss.} \textbf{13}, 711 (1925).

\item\label{ref:1.3} C. J. Davisson and L. H. Germer, \emph{Phys. Rev.} \textbf{30}, 705
(1927).

\item\label{ref:1.4} W. Heisenberg, \emph{A. Phys.} \textbf{33}, 879 (1925); M. Born and P.
Jordan, \emph{Z. f. Phys.} \textbf{34}, 858 (1925); P. A. M. Dirac, \emph{Proc.
Roy. Soc.} A\textbf{109}, 642 (1925); M. Born, W. Heisenberg, and P. Jordan,
\emph{Z. f. Phys.} \textbf{35}, 557 (1925); W. Pauli, \emph{Z. f. Phys.}
\textbf{36}, 336 (1926). These papers are reprinted in \emph{Sources of
Quantum Mechanics}, ed. by B. L. van der Waerden (Dover Publications, Inc.,
New York, 1968).

\item\label{ref:1.5} E. Schr\"odinger, \emph{Ann. Phys.} \textbf{79}, 361, 489; \textbf{80},
437; \textbf{81}, 109 (1926). These papers are reprinted in English,
unfortunately in a somewhat abridged form, in \emph{Wave Mechanics}, Ref.~\hyperref[ref:1.1]{1}.
Also see \emph{Collected Papers on Wave Mechanics}, trans. by J. F. Schearer
and W. M. Deans (Blackie and Son, London, 1928).

\item\label{ref:1.6} See, e.g., P. A. M. Dirac, \emph{The Development of Quantum Theory}
(Gordon and Breach, New York, 1971). Also see Dirac's obituary of Schr\"odinger,
\emph{Nature} \textbf{189}, 355 (1961), and his article in \emph{Scientific
American} \textbf{208}, 45 (1963).

\item\label{ref:1.7} O. Klein, \emph{Z. f. Phys.} \textbf{37}, 895 (1926). Also see V. Fock,
\emph{Z. f. Phys.} \textbf{38}, 242 (1926); \emph{ibid.} \textbf{39}, 226
(1926).

\item\label{ref:1.8} W. Gordon, \emph{Z. f. Phys.} \textbf{40}, 117 (1926).

\item\label{ref:1.9} For the details of the calculation, see, e.g., L. I. Schiff,
\emph{Quantum Mechanics}, 3rd edn, (McGraw-Hill, Inc. New York, 1968): Section
51.

\item\label{ref:1.10} F. Paschen, \emph{Ann. Phys.} \textbf{50}, 901 (1916). These
experiments were actually carried out using He$^+$ because its fine structure
splitting is 16 times larger than for hydrogen. The fine structure of spectral
lines was first discovered interferometrically by A. A. Michelson, \emph{Phil.
Mag.} \textbf{31}, 338 (1891); \emph{ibid.}, \textbf{34}, 280 (1892).

\item[10a.]\phantomsection\label{ref:1.10a}\label{ref:1.10c} A. Sommerfeld, \emph{M\"unchner Berichte} 1915, pp. 425, 429;
\emph{Ann. Phys.} \textbf{51}, 1, 125 (1916). Also see W. Wilson, \emph{Phil.
Mag.} \textbf{29}, 795 (1915).

\item\label{ref:1.11} G. E. Uhlenbeck and S. Goudsmit, \emph{Naturwiss.} \textbf{13}, 953
(1925); \emph{Nature} \textbf{117}, 264 (1926). The electron spin had been
earlier suggested for other reasons by A. H. Compton, \emph{J. Frank. Inst.}
\textbf{192}, 145 (1921).

\item\label{ref:1.12} The general formula for Zeeman splitting in one-electron atoms had been
discovered empirically by A. Land\'e, \emph{Z. f. Phys.} \textbf{5}, 231
(1921); \emph{ibid.}, \textbf{7}, 398 (1921); \emph{ibid.}, \textbf{15}, 189
(1923); \emph{ibid.}, \textbf{19}, 112 (1923). At the time, the extra
non-orbital angular momentum appearing in this formula was thought to be the
angular momentum of the atomic `core;' A. Sommerfeld, \emph{Ann. Phys.}
\textbf{63}, 221 (1920); \emph{ibid.}, \textbf{70}, 32 (1923). It was only
later that the extra angular momentum was recognized, as in Ref.~\hyperref[ref:1.11]{11}, to be due
to the spin of the electron.

\item\label{ref:1.13} W. Heisenberg and P. Jordan, \emph{Z. f. Phys.} \textbf{37}, 263
(1926); C. G. Darwin, \emph{Proc. Roy. Soc.} A\textbf{116}, 227 (1927).
Darwin says that several authors did this work at about the same time, while
Dirac quotes only Darwin.

\item\label{ref:1.14} L. H. Thomas, \emph{Nature} \textbf{117}, 514 (1926). Also see S.
Weinberg, \emph{Gravitation and Cosmology}, (Wiley, New York, 1972): Section
5.1.

\item\label{ref:1.15} P. A. M. Dirac, \emph{Proc. Roy. Soc.} A\textbf{117}, 610 (1928). Also
see Dirac, \emph{ibid.}, A\textbf{118}, 351 (1928), for the application of this
theory to the calculation of the Zeeman and Paschen--Back effects and the
relative strengths of lines within fine-structure multiplets.

\item\label{ref:1.16} For the probabilistic interpretation of non-relativistic quantum
mechanics, see M. Born \emph{Z. f. Phys.} \textbf{37}, 863 (1926); \emph{ibid.},
\textbf{38}, 803 (1926) (reprinted in an abridged English version in \emph{Wave
Mechanics}, Ref.~\hyperref[ref:1.41]{41}); G. Wentzel, \emph{Z. f. Phys.} \textbf{40}, 590 (1926);
W. Heisenberg \emph{Z. f. Phys.} \textbf{43}, 172 (1927). N. Bohr.
\emph{Nature} \textbf{121}, 580 (1928); \emph{Naturwissenschaften} \textbf{17},
483 (1929); \emph{Electrons et Photons---Rapports et Discussions du V$^e$
Conseil de Physique Solvay} (Gauthier-Villars, Paris, 1928).

\item\label{ref:1.17} Conversation between Dirac and J. Mehra, March 28, 1969, quoted by
Mehra in \emph{Aspects of Quantum Theory}, ed. by A. Salam and E. P. Wigner
(Cambridge University Press, Cambridge, 1972).

\item\label{ref:1.18} G. Gamow, \emph{Thirty Years that Shook Physics}, (Doubleday and Co.,
Garden City, NY, 1966): p. 125.

\item\label{ref:1.19} W. Pauli, \emph{Z. f. Phys.} \textbf{37}, 263 (1926); \textbf{43}, 601
(1927).

\item\label{ref:1.20} C. G. Darwin, \emph{Proc. Roy. Soc.} A\textbf{118}, 654 (1928);
\emph{ibid.}, A\textbf{120}, 621 (1928).

\item\label{ref:1.21} W. Gordon, \emph{Z. f. Phys.} \textbf{48}, 11 (1928).

\item\label{ref:1.22} P. A. M. Dirac, \emph{Proc. Roy. Soc.} A\textbf{126}, 360 (1930); also
see Ref.~\hyperref[ref:1.47]{47}.

\item\label{ref:1.23} E. C. Stoner, \emph{Phil. Mag.} \textbf{48}, 719 (1924).

\item\label{ref:1.24} W. Pauli, \emph{Z. f. Phys.} \textbf{31}, 765 (1925).

\item\label{ref:1.25} W. Heisenberg, \emph{Z. f. Phys.} \textbf{38}, 411 (1926); \emph{ibid.},
\textbf{39}, 499 (1926); P. A. M. Dirac, \emph{Proc. Roy. Soc.}
A\textbf{112}, 661 (1926); W. Pauli, \emph{Z. f. Phys.} \textbf{41}, 81
(1927); J. C. Slater, \emph{Phys. Rev.} \textbf{34}, 1293 (1929).

\item\label{ref:1.26} E. Fermi, \emph{Z. f. Phys.} \textbf{36}, 902 (1926); \emph{Rend.
Accad. Lincei} \textbf{3}, 145 (1926).

\item\label{ref:1.27} P. A. M. Dirac, Ref.~\hyperref[ref:1.25]{25}.

\item[27a.]\phantomsection\label{ref:1.27a} P. A. M. Dirac, First W. R. Crane Lecture at the University of
Michigan, April 17, 1989, unpublished.

\item\label{ref:1.28} H. Weyl, \emph{The Theory of Groups and Quantum Mechanics}, translated
from the second (1931) German edition by H. P. Robertson (Dover Publications,
Inc., New York): Chapter IV, Section 12. Also see P. A. M. Dirac, \emph{Proc.
Roy. Soc.} A\textbf{133}, 61 (1931).

\item\label{ref:1.29} J. R. Oppenheimer, \emph{Phys. Rev.} \textbf{35}, 562 (1930); I. Tamm,
\emph{Z. f. Phys.} \textbf{62}, 545 (1930).

\item[29a.]\phantomsection\label{ref:1.29a} P. A. M. Dirac, \emph{Proc. Roy. Soc.} \textbf{133}, 60 (1931).

\item\label{ref:1.30} C. D. Anderson, \emph{Science} \textbf{76}, 238 (1932); \emph{Phys.
Rev.} \textbf{43}, 491 (1933). The latter paper is reprinted in \emph{Foundations
of Nuclear Physics}, ed. by R. T. Beyer (Dover Publications, Inc., New York,
1949).

\item[30a.]\phantomsection\label{ref:1.30a} J. Schwinger, `A Report on Quantum Electrodynamics,' in \emph{The
Physicist's Conception of Nature} (Reidel, Dordrecht, 1973): p.415.

\item\label{ref:1.31} W. Pauli, \emph{Handbuch der Physik} (Julius Springer, Berlin,
1932--1933); \emph{Rev. Mod. Phys.} \textbf{13}, 203 (1941).

\item\label{ref:1.32} Born, Heisenberg, and Jordan, Ref.~\hyperref[ref:1.4]{4}, Section 3.

\item[32a.]\phantomsection\label{ref:1.32a} P. Ehrenfest, \emph{Phys. Z.} \textbf{7}, 528 (1906).

\item\label{ref:1.33} Born and Jordan, Ref.~\hyperref[ref:1.4]{4}. Unfortunately the relevant parts of this paper
are not included in the reprint collection \emph{Sources of Quantum Mechanics},
cited in Ref.~\hyperref[ref:1.4]{4}.

\item\label{ref:1.34} P. A. M. Dirac, \emph{Proc. Roy. Soc.} A\textbf{112}, 661 (1926):
Section 5. For a more accessible derivation, see L. I. Schiff, \emph{Quantum
Mechanics}, 3rd edn. (McGraw-Hill Book Company, New York, 1968): Section 44.

\item[34a.]\phantomsection\label{ref:1.34a} A. Einstein, \emph{Phys. Z.} \textbf{18}, 121 (1917); reprinted
in English in van der Waerden, Ref.~\hyperref[ref:1.4]{4}.

\item\label{ref:1.35} P. A. M. Dirac, \emph{Proc. Roy. Soc.} A\textbf{114}, 243 (1927);
reprinted in \emph{Quantum Electrodynamics}, ed. by J. Schwinger (Dover
Publications, Inc., New York, 1958).

\item\label{ref:1.36} P. A. M. Dirac, \emph{Proc. Roy. Soc.} A\textbf{114}, 710 (1927).

\item[36a.]\phantomsection\label{ref:1.36a} V. F. Weisskopf and E. Wigner, \emph{Z. f. Phys.} \textbf{63}, 54
(1930).

\item[36b.]\phantomsection\label{ref:1.36b} E. Fermi, \emph{Lincei Rend.} \textbf{9}, 881 (1929); \textbf{12},
431 (1930); \emph{Rev. Mod. Phys.} \textbf{4}, 87 (1932).

\item\label{ref:1.37} P. Jordan and W. Pauli, \emph{Z. f. Phys.} \textbf{47}, 151 (1928).

\item\label{ref:1.38} N. Bohr and L. Rosenfeld, \emph{Kon. dansk. vid. Selsk., Mat.-Fys.
Medd.} \textbf{XII}, No. 8 (1933) (translation in \emph{Selected Papers of Leon
Rosenfeld}, ed. by R. S. Cohen and J. Stachel (Reidel, Dordrecht, 1979));
\emph{Phys. Rev.} \textbf{78}, 794 (1950).

\item\label{ref:1.39} P. Jordan, \emph{Z. f. Phys.} \textbf{44}, 473 (1927). Also see P.
Jordan and O. Klein, \emph{Z. f. Phys.} \textbf{45}, 751 (1929); P. Jordan,
\emph{Phys. Zeit.} \textbf{30}, 700 (1929).

\item\label{ref:1.40} P. Jordan and E. Wigner, \emph{Z. f. Phys.} \textbf{47}, 631 (1928).
This article is reprinted in \emph{Quantum Electrodynamics}, Ref.~\hyperref[ref:1.35]{35}.

\item[40a.]\phantomsection\label{ref:1.40a} M. Fierz, \emph{Helv. Phys. Acta} \textbf{12}, (1939); W. Pauli,
\emph{Phys. Rev.} \textbf{58}, 716 (1940); W. Pauli and F. J. Belinfante,
\emph{Physica} \textbf{7}, 177 (1940).

\item\label{ref:1.41} W. Heisenberg and W. Pauli, \emph{Z. f. Phys.} \textbf{56}, 1 (1929);
\emph{ibid.}, \textbf{59}, 168 (1930).

\item\label{ref:1.42} P. A. M. Dirac, \emph{Proc. Roy. Soc.} A\textbf{136}, 453 (1932); P.
A. M. Dirac, V. A. Fock, and B. Podolsky, \emph{Phys. Zeit. der Sowjetunion}
\textbf{2}, 468 (1932); P. A. M. Dirac, \emph{Phys. Zeit. der Sowjetunion}
\textbf{3}, 64 (1933). The latter two articles are reprinted in \emph{Quantum
Electrodynamics}, Ref.~\hyperref[ref:1.35]{35}, pp. 29 and 312. Also see L. Rosenfeld, \emph{Z. f.
Phys.} \textbf{76}, 729 (1932).

\item[42a.]\phantomsection\label{ref:1.42a} P. A. M. Dirac, \emph{Proc. Roy. Soc. London} A\textbf{136}, 453
(1932).

\item\label{ref:1.43} E. Fermi, \emph{Z. f. Phys.} \textbf{88}, 161 (1934). Fermi quotes
unpublished work of Pauli for the proposition that an unobserved neutral
particle is emitted along with the electron in beta decay. This particle was
called the neutrino to distinguish it from the recently discovered neutron.

\item[43a.]\phantomsection\label{ref:1.43a} V. Fock, \emph{C. R. Leningrad} 1933, p. 267.

\item\label{ref:1.44} W. H. Furry and J. R. Oppenheimer, \emph{Phys. Rev.} \textbf{45}, 245
(1934). This paper uses a density matrix formalism developed by P. A. M.
Dirac, \emph{Proc. Camb. Phil. Soc.} \textbf{30}, 150 (1934). Also see R. E.
Peierls, \emph{Proc. Roy. Soc.} \textbf{146}, 420 (1934); W. Heisenberg,
\emph{Z. f. Phys.} \textbf{90}, 209 (1934); L. Rosenfeld, \emph{Z. f. Phys.}
\textbf{76}, 729 (1932).

\item\label{ref:1.45} W. Pauli and V. Weisskopf, \emph{Helv. Phys. Acta} \textbf{7}, 709
(1934), reprinted in English translation in A. I. Miller, \emph{Early Quantum
Electrodynamics} (Cambridge University Press, Cambridge, 1994). Also see W.
Pauli, \emph{Ann. Inst. Henri Poincar\'e} \textbf{6}, 137 (1936).

\item\label{ref:1.46} O. Klein and Y. Nishina, \emph{Z. f. Phys.} \textbf{52}, 853 (1929);
Y. Nishina, \emph{ibid.}, 869 (1929); also see I. Tamm, \emph{Z. f. Phys.}
\textbf{62}, 545 (1930).

\item\label{ref:1.47} P. A. M. Dirac, \emph{Proc. Camb. Phil. Soc.} \textbf{26}, 361 (1930).

\item\label{ref:1.48} C. M\o ller, \emph{Ann. d. Phys.} \textbf{14}, 531, 568 (1932).

\item\label{ref:1.49} H. Bethe and W. Heitler, \emph{Proc. Roy. Soc.} A\textbf{146}, 83
(1934); also see G. Racah, \emph{Nuovo Cimento} \textbf{11}, No. 7 (1934);
\emph{ibid.}, \textbf{13}, 69 (1936).

\item\label{ref:1.50} H. J. Bhabha, \emph{Proc. Roy. Soc.} A\textbf{154}, 195 (1936).

\item[50a.]\phantomsection\label{ref:1.50a} J. F. Carlson and J. R. Oppenheimer, \emph{Phys. Rev.}
\textbf{51}, 220 (1937).

\item\label{ref:1.51} P. Ehrenfest and J. R. Oppenheimer, \emph{Phys. Rev.} \textbf{37}, 333
(1931).

\item\label{ref:1.52} W. Heitler and G. Herzberg, \emph{Naturwiss.} \textbf{17}, 673 (1929);
F. Rasetti, \emph{Z. f. Phys.} \textbf{61}, 598 (1930).

\item\label{ref:1.53} J. Chadwick, \emph{Proc. Roy. Soc.} A\textbf{136}, 692 (1932). This
article is reprinted in \emph{The Foundations of Nuclear Physics}, Ref.~\hyperref[ref:1.30]{30}.

\item\label{ref:1.54} W. Heisenberg, \emph{Z. f. Phys.} \textbf{77}, 1 (1932); also see I.
Curie-Joliot and F. Joliot, \emph{Compt. Rend.} \textbf{194}, 273 (1932).

\item[54a.]\phantomsection\label{ref:1.54a} For references, see L. M. Brown and H. Rechenberg, \emph{Hist.
Stud. in Phys. and Bio. Science}, \textbf{25}, 1 (1994).

\item\label{ref:1.55} H. Yukawa, \emph{Proc. Phys.-Math. Soc. (Japan)} (3) \textbf{17}, 48
(1935). This article is reprinted in \emph{The Foundations of Nuclear Physics},
Ref.~\hyperref[ref:1.30]{30}.

\item\label{ref:1.56} S. H. Neddermeyer and C. D. Anderson, \emph{Phys. Rev.} \textbf{51},
884 (1937); J. C. Street and E. C. Stevenson, \emph{Phys. Rev.} \textbf{52},
1003 (1937).

\item[56a.]\phantomsection\label{ref:1.56a} L. Nordheim and N. Webb, \emph{Phys. Rev.} \textbf{56}, 494
(1939).

\item\label{ref:1.57} M. Conversi, E. Pancini, and O. Piccioni, \emph{Phys. Rev.} \textbf{71},
209L (1947).

\item\label{ref:1.58} S. Sakata and T. Inoue, \emph{Prog. Theor. Phys.} \textbf{1}, 143
(1946); R. E. Marshak and H. A. Bethe, \emph{Phys. Rev.} \textbf{77}, 506
(1947).

\item\label{ref:1.59} C. M. G. Lattes, G. P. S. Occhialini, and C. F. Powell, \emph{Nature}
\textbf{160}, 453, 486 (1947).

\item\label{ref:1.60} G. D. Rochester and C. C. Butler, \emph{Nature} \textbf{160}, 855
(1947).

\item\label{ref:1.61} J. R. Oppenheimer, \emph{Phys. Rev.} \textbf{35}, 461 (1930).

\item\label{ref:1.62} I. Waller, \emph{Z. f. Phys.} \textbf{59}, 168 (1930); \emph{ibid.},
\textbf{61}, 721, 837 (1930); \emph{ibid.}, \textbf{62}, 673 (1930).

\item\label{ref:1.63} V. F. Weisskopf, \emph{Z. f. Phys.} \textbf{89}, 27 (1934), reprinted
in English translation in \emph{Early Quantum Electrodynamics}, Ref.~\hyperref[ref:1.45]{45};
\emph{ibid.}, \textbf{90}, 817 (1934). The electromagnetic self energy is
calculated only to lowest order in $\alpha$ in these references; the proof that
the divergence is only logarithmic in all orders of perturbation theory was
given by Weisskopf; \emph{Phys. Rev.} \textbf{56}, 72 (1939). (This last
article is reprinted in \emph{Quantum Electrodynamics}, Ref.~\hyperref[ref:1.35]{35}).

\item\label{ref:1.64} P. A. M. Dirac, XVII Conseil Solvay de Physique, p. 203 (1933),
reprinted in \emph{Early Quantum Electrodynamics}, Ref.~\hyperref[ref:1.45]{45}. For subsequent
calculations based on less restrictive assumptions, see W. Heisenberg,
\emph{Z. f. Phys.} \textbf{90}, 209 (1934); \emph{Sachs. Akad. Wiss.}
\textbf{86}, 317 (1934); R. Serber, \emph{Phys. Rev.} \textbf{43}, 49 (1935);
E. A. Uehling, \emph{Phys. Rev.} \textbf{48}, 55 (1935); W. Pauli and M. Rose,
\emph{Phys. Rev.} \textbf{49}, 462 (1936). Also see Furry and Oppenheimer,
Ref.~\hyperref[ref:1.44]{44}; Peierls, Ref.~\hyperref[ref:1.44]{44}; Weisskopf, Ref.~\hyperref[ref:1.63]{63}.

\item\label{ref:1.65} H. Euler and B. Kockel, \emph{Naturwiss.} \textbf{23}, 246 (1935); W.
Heisenberg and H. Euler, \emph{Z. f. Phys.} \textbf{98}, 714 (1936).

\item\label{ref:1.66} P. A. M. Dirac, \emph{Proc. Camb. Phil. Soc.} \textbf{30}, 150 (1934).

\item\label{ref:1.67} W. Heisenberg, \emph{Z. f. Phys.} \textbf{90}, 209 (1934).

\item\label{ref:1.68} N. Kemmer and V. F. Weisskopf. \emph{Nature} \textbf{137}, 659 (1936).

\item[68a.]\phantomsection\label{ref:1.68a} F. Bloch and A. Nordsieck, \emph{Phys. Rev.} \textbf{52}, 54
(1937). Also see W. Pauli and M. Fierz, \emph{Nuovo Cimento} \textbf{15}, 167
(1938), reprinted in English translation in \emph{Early Quantum
Electrodynamics}, Ref.~\hyperref[ref:1.45]{45}.

\item\label{ref:1.69} S. M. Dancoff, \emph{Phys. Rev.} \textbf{55}, 959 (1939).

\item[69a.]\phantomsection\label{ref:1.69a} H. W. Lewis, \emph{Phys. Rev.} \textbf{73}, 173 (1948); S.
Epstein, \emph{Phys. Rev.} \textbf{73}, 177 (1948). Also see J. Schwinger,
Ref.~\hyperref[ref:1.84]{84}; Z. Koba and S. Tomonaga, \emph{Prog. Theor. Phys.} \textbf{3/3}, 290
(1948).

\item[69b.]\phantomsection\label{ref:1.69b} J. Schwinger, in \emph{The Birth of Particle Physics}, ed. by L.
Brown and L. Hoddeson (Cambridge University Press, Cambridge, 1983): p. 336.

\item\label{ref:1.70} W. Heisenberg, \emph{Ann. d. Phys.} \textbf{32}, 20 (1938), reprinted
in English translation in \emph{Early Quantum Electrodynamics}, Ref.~\hyperref[ref:1.45]{45}.

\item[70a.]\phantomsection\label{ref:1.70a} G. Wentzel, \emph{Z. f. Phys.} \textbf{86}, 479, 635 (1933);
\emph{Z. f. Phys.} \textbf{87}, 726 (1934); M. Born and L. Infeld, \emph{Proc.
Roy. Soc.} A\textbf{150}, 141 (1935); W. Pauli, \emph{Ann. Inst. Henri
Poincar\'e} \textbf{6}, 137 (1936).

\item\label{ref:1.71} J. A. Wheeler, \emph{Phys. Rev.} \textbf{52}, 1107 (1937).

\item\label{ref:1.72} W. Heisenberg, \emph{Z. f. Phys.} \textbf{120}, 513, 673 (1943);
\emph{Z. Naturforsch.} \textbf{1}, 608 (1946). Also see C. M\o ller, \emph{Kon.
Dansk. Vid. Sels. Mat.-Fys. Medd.} \textbf{23}, No. 1 (1945); \emph{ibid.}
\textbf{23}, No. 19, (1946).

\item\label{ref:1.73} See, e.g., G. Chew, \emph{The S-Matrix Theory of Strong Interactions}
(W. A. Benjamin, Inc. New York, 1961).

\item\label{ref:1.74} J. A. Wheeler and R. P. Feynman, \emph{Rev. Mod. Phys.} \textbf{17},
157 (1945), \emph{ibid.}, \textbf{21}, 425 (1949). For further references and a
discussion of the application of action-at-a-distance theories in cosmology,
see S. Weinberg \emph{Gravitation and Cosmology}, (Wiley, 1972): Section 16.3.

\item\label{ref:1.75} P. A. M. Dirac, \emph{Proc. Roy. Soc.} A\textbf{180}, 1 (1942). For a
criticism, see W. Pauli, \emph{Rev. Mod. Phys.} \textbf{15}, 175 (1943). For a
review of classical theories of this type, and of yet other attempts to solve
the problem of infinities, see R. E. Peierls in \emph{Rapports du 8$^{\rm me}$
Conseil de Physique Solvay 1948} (R. Stoops, Brussels, 1950): p. 241.

\item\label{ref:1.76} V. F. Weisskopf, \emph{Kon. Dan. Vid. Sel., Mat.-fys. Medd.}
\textbf{XIV}, No. 6 (1936), especially p. 34 and pp. 5--6. This article is
reprinted in \emph{Quantum Electrodynamics}, Ref.~\hyperref[ref:1.35]{35}, and in English
translation in \emph{Early Quantum Electrodynamics}, Ref.~\hyperref[ref:1.45]{45}. Also see W.
Pauli and M. Fierz, Ref.~\hyperref[ref:1.68a]{68a}; H. A. Kramers, Ref.~\hyperref[ref:1.79a]{79a}.

\item\label{ref:1.77} S. Pasternack, \emph{Phys. Rev.} \textbf{54}, 1113 (1938). This
suggestion was based on experiments of W.V. Houston, \emph{Phys. Rev.}
\textbf{51}, 446 (1937); R. C. Williams, \emph{Phys. Rev.} \textbf{54}, 558
(1938). For a report of contrary data, see J. W. Drinkwater, O. Richardson,
and W. E. Williams, \emph{Proc. Roy. Soc.} \textbf{174}, 164 (1940).

\item\label{ref:1.78} E. A. Uehling, Ref.~\hyperref[ref:1.64]{64}.

\item\label{ref:1.79} W. E. Lamb, Jr and R. C. Retherford, \emph{Phys. Rev.} \textbf{72}, 241
(1947). This article is reprinted in \emph{Quantum Electrodynamics}, Ref.~\hyperref[ref:1.35]{35}.

\item[79a.]\phantomsection\label{ref:1.79a} H. A. Kramers, \emph{Nuovo Cimento} \textbf{15}, 108 (1938),
reprinted in English translation in \emph{Early Quantum Electrodynamics}, Ref.
45; \emph{Ned. T. Natuurk.} \textbf{11}, 134 (1944); \emph{Rapports du 8$^{\rm
me}$ Conseil de Physique Solvay 1948} (R. Stoops, Brussels, 1950).

\item\label{ref:1.80} H. A. Bethe, \emph{Phys. Rev.} \textbf{72}, 339 (1947). This article is
reprinted in \emph{Quantum Electrodynamics}, Ref.~\hyperref[ref:1.35]{35}.

\item\label{ref:1.81} J. B. French and V. F. Weisskopf; \emph{Phys. Rev.} \textbf{75}, 1240
(1949); N. M. Kroll and W. E. Lamb, \emph{ibid.}, \textbf{75}, 388 (1949); J.
Schwinger, \emph{Phys. Rev.} \textbf{75}, 898 (1949); R. P. Feynman,
\emph{Rev. Mod. Phys.} \textbf{20}, 367 (1948); \emph{Phys. Rev.}, \textbf{74},
939, 1430 (1948); \textbf{76}, 749, 769 (1949); \textbf{80}, 440 (1950); H.
Fukuda, Y. Miyamoto, and S. Tomonaga, \emph{Prog. Theor. Phys. Rev. Mod.
Phys.} \textbf{4}, 47, 121 (1948). The article by Kroll and Lamb is reprinted
in \emph{Quantum Electrodynamics}, Ref.~\hyperref[ref:1.35]{35}.

\item\label{ref:1.82} J. E. Nafe, E. B. Nelson, and I. I. Rabi, \emph{Phys. Rev.} \textbf{71},
914 (1947); D. E. Nagel, R. S. Julian, and J. R. Zacharias, \emph{Phys. Rev.}
\textbf{72}, 973 (1947).

\item\label{ref:1.83} P. Kusch and H. M. Foley, \emph{Phys. Rev.} \textbf{72}, 1256 (1947).

\item[83a.]\phantomsection\label{ref:1.83a} G. Breit, \emph{Phys. Rev.} \textbf{71}, 984 (1947). Schwinger in
Ref.~\hyperref[ref:1.84]{84} includes a corrected version of Breit's results.

\item\label{ref:1.84} J. Schwinger, \emph{Phys. Rev.} \textbf{73}, 416 (1948). This article
is reprinted in \emph{Quantum Electrodynamics}, Ref.~\hyperref[ref:1.35]{35}.

\item\label{ref:1.85} J. Schwinger, \emph{Phys. Rev.} \textbf{74}, 1439 (1948); \emph{ibid.},
\textbf{75}, 651 (1949); \emph{ibid.}, \textbf{76}, 790 (1949); \emph{ibid.},
\textbf{82}, 664, 914 (1951); \emph{ibid.}, \textbf{91}, 713 (1953); \emph{Proc.
Nat. Acad. Sci.} \textbf{37}, 452 (1951). All but the first two of these
articles are reprinted in \emph{Quantum Electrodynamics}, Ref.~\hyperref[ref:1.35]{35}.

\item\label{ref:1.86} S. Tomonaga, \emph{Prog. Theor. Phys. Rev. Mod. Phys.} \textbf{1}, 27
(1946); Z. Koba, T. Tati, and S. Tomonaga, \emph{ibid.}, \textbf{2}, 101
(1947); S. Kanesawa and S. Tomonaga, \emph{ibid.}, \textbf{3}, 1, 101 (1948);
S. Tomonaga, \emph{Phys. Rev.} \textbf{74}, 224 (1948); D. Ito, Z. Koba, and S.
Tomonaga, \emph{Prog. Theor. Phys.} \textbf{3}, 276 (1948); Z. Koba and S.
Tomonaga, \emph{ibid.}, \textbf{3}, 290 (1948). The first and fourth of these
articles are reprinted in \emph{Quantum Electrodynamics}, Ref.~\hyperref[ref:1.35]{35}.

\item\label{ref:1.87} R. P. Feynman, \emph{Rev. Mod. Phys.} \textbf{20}, 367 (1948);
\emph{Phys. Rev.} \textbf{74}, 939, 1430 (1948); \emph{ibid.}, \textbf{76},
749, 769 (1949); \emph{ibid.} \textbf{80}, 440 (1950). All but the second and
third of these articles are reprinted in \emph{Quantum Electrodynamics}, Ref.
35.

\item\label{ref:1.88} F. J. Dyson, \emph{Phys. Rev.} \textbf{75}, 486, 1736 (1949). These
articles are reprinted in \emph{Quantum Electrodynamics}, Ref.~\hyperref[ref:1.35]{35}.

\item[88a.]\phantomsection\label{ref:1.88a} H. Fr\"ohlich, W. Heitler, and B. Kahn, \emph{Proc. Roy. Soc.}
A\textbf{171}, 269 (1939); \emph{Phys. Rev.} \textbf{56}, 961 (1939).

\item[88b.]\phantomsection\label{ref:1.88b} W. E. Lamb, Jr, \emph{Phys. Rev.} \textbf{56}, 384 (1939);
\emph{Phys. Rev.} \textbf{57}, 458 (1940).

\item\label{ref:1.89} Quoted by R. Serber, in \emph{The Birth of Particle Physics}, Ref.~\hyperref[ref:1.69b]{69b},
p. 270.
\end{enumerate}
